\documentclass[10pt]{article} 
\renewcommand{\thempfootnote}{\arabic{mpfootnote}}

\usepackage[hidelinks]{hyperref}

\usepackage[utf8]{inputenc}

\usepackage[T1]{fontenc}
\usepackage[ttdefault=true]{AnonymousPro}

\usepackage{multicol}
\setlength{\columnsep}{1cm}

\usepackage{abc}
\usepackage{graphicx}
\usepackage{musicography}

\usepackage[margin=0.5in]{geometry} 
\geometry{letterpaper}

\usepackage{fancyhdr} 
\pagestyle{fancy} % options: empty , plain , fancy
\renewcommand{\headrulewidth}{0pt} % customise the layout...
\lhead{}\chead{}\rhead{}
\lfoot{}\cfoot{\thepage}\rfoot{}

\begin{document}

\thispagestyle{empty}

\begin{abc}[name=lowerys,options={-O= +c -F single.fmt}]
X:1
T:Lowery's Quadrille
R:Fiddle tuned ADae
Q:1/2=110
L:1/4
M:C|
S:John Summers to Armin Barnett to the Canote Brothers
K:D
|:"D"A,DFA | "G"B/d/B "D"A2 | dBAF | "A"E/D/F/E/ D/B,/-B,| 
"D"A,DFA | "G"B/d/B "D"A2 | dBAF | "A"E/F/E "D"D2 :|
|:"D"D/F/A/c/ d/A/-A | "G"B-B/A/ Bd/B/ | "D"AA/B/ A/F/-F |  "A"E-E/F/ EF/E/ | 
"D"D/F/A/c/ d/A/-A | "G"B-B/A/ Bd/B/ |1 "D"A/F/-F/D/ "A"E/F/E | "D"D-D/E/ DF/E/ :|2"A"A/A/c/c/ e/e/g/g/ | f/d/e/c/ "D"de | 
|:"D"f-f/a/ f/e/d/d/ | "G"e/f/e/d/ B/d/A | "D"f/a/-a/f/ aa/b/ | "A"afed | 
"D"f-f/a/ f/e/d/d/ | "G"e/f/e/d/ B/d/ (3A/B/c/ | "D"d!+!AF!+!D | "A"A,/D/-D "D"!3!!0![DD]2 :| 
\end{abc}

\vskip 0.5in

\noindent\makebox[\textwidth][c]{
  \begin{minipage}{7.0in}
    \setlength{\parindent}{1.5em}

The tune \textbf{Midnight on the Water} is attributed to the Texan
fiddler Luke\footnote{I also found `Lucas' and `Lewis Alexander':
  \\ \url{https://utaholdtimefiddlers.org/author/fiddler/page/2/}
  \\ \url{https://en.wikipedia.org/wiki/Benny_Thomasson}} ``Red Luke''
Thomasson, dating back to around
1900.\footnote{\url{https://tunearch.org/wiki/Annotation:Midnight_on_the_Water}}
It gained notable recognition after being recorded by his son, fiddler
Benny Thomasson (1909 – 1984), in 1970 for his album `Country Fiddling
from the Big State,' subsequently becoming popular within Old Time
music circles. There is speculation that certain sections of
\textbf{Midnight on the Water} bear resemblance to \textbf{Old
  Paint},\footnote{I think `Old Paint' refers to a horse, but I can't
  find confirmation. \textit{[``Paint's a good pony, and she paces
      when she can'' is a line in the chorus, at least as sung by
      Bruce Molsky. ---Josh]}} a tune by Jess Morris, potentially
indicating mutual inspiration, given the proximity of the Morrises and
Thomassons in the Texas panhandle during that time period. Often
performed in Deadman's Tuning (DDAD) alongside the standard GDAE
tuning, both are viable options, yet Deadman's Tuning lends the
composition a deep resonance, enriching its serene ebbing character.

---Kitty Steetle

  \end{minipage}
}

\end{document}
